\section{Related work and contribution boundary}\label{sec:related}
\paragraph{Engineering knowledge and executable contracts.}
Process knowledge representation links decision context, alternatives and dependencies in complex engineered systems \cite{wang2021process}. Decision knowledge graphs integrate and navigate engineering decisions \cite{hao2021knowledge}; manufacturing work further uses graph learning for knowledge reasoning \cite{su2025manufacturing}. Our representation is a typed dependency graph and a monitored state machine. It performs deterministic applicability checks, selection and invalidation. It does not learn relations or claim general semantic reasoning. Contract theory provides a much broader algebra of assumptions, guarantees and composition \cite{benveniste2018contracts}. We implement a specific operating contract whose statistical part concerns a fixed policy and an explicitly defined localization event.

\paragraph{Computational decisions in autonomous systems.}
Metareasoning and anytime algorithms treat computation as a decision with opportunity cost \cite{russell1991meta,zilberstein1996anytime,hay2012}. Pant et al. already use contract-driven anytime estimation and control with robust model-predictive control and a physical aerial platform \cite{pant2021}. Their work is the closest precedent for connecting computational contracts to autonomy. We do not claim the first such connection, robust closed-loop control or a new general optimal-stopping theorem. Perception scheduling can preserve stability while trading latency against precision \cite{aldana2023}; computational abstraction relates robot task accuracy to compute expenditure \cite{ghosh2021}; kernel studies characterize practical computational demand \cite{sifat2023}. Here the distinctive execution object is a correlated source-version stream, and the evidence object is an event-specific, population-level risk card that must retain that stream's policy and installation identity.

\paragraph{Registration and state consistency.}
ICP variants and point-cloud registration are established methods \cite{besl1992icp,chen1992,segal2009gicp,pomerleau2015}. Our contribution lies in exposing a retained solve and governing its publication, not in replacing the registration objective. ICP curvature is also an imperfect surrogate for error uncertainty \cite{censi2007,bonnabel2016,brossard2020}. We therefore retain the joint robot/map state, propagate its covariance, and obtain empirical risk bounds from entire executed episodes rather than interpreting a Hessian as calibrated confidence. Delayed-measurement estimation already addresses temporal consistency \cite{barshalom2002}; our immutable-source replacement handles correlated revisions of one solve rather than fusing them as independent observations. Complete localization systems such as KISS-ICP and FAST-LIO2 address broader odometry and mapping problems \cite{kiss2023,xu2022}. They are not substituted for same-task scheduling controls: our task uses a supplied structural prior and a common scan registration interface.

\paragraph{Risk control and runtime assurance.}
Exact binomial inversion supplies the one-sided confidence bounds \cite{clopper1934}. Learn-then-test separates fitted predictive algorithms from risk-controlled selection \cite{angelopoulos2025ltt}; conformal risk control and conformal policy control provide other principled control constructions \cite{angelopoulos2024crc,prinster2026}. We use a finite, pre-frozen catalogue with Bonferroni allocation because the events are binary and simultaneous bounds directly support the requested selector. This is not a new statistical bound. Runtime-assurance architectures can switch to a safety controller \cite{mehmood2022simplex}; our refusal withdraws statistical attachment or retains an uncorrected pose and supplies no physical safety guarantee. The engineering contribution is the executable correspondence between requirements, numerical execution and evidence lifetime, assessed against an independently coded procedure that can enforce the same semantics.

\begin{table}[t]\centering\footnotesize
\caption{Closest mechanism precedents and the object evaluated here. The rows describe supported research objects and do not assert that a precedent lacks every other capability.}\label{tab:mechanism-precedents}
\begin{tabularx}{\linewidth}{@{}p{.20\linewidth}XXX@{}}\toprule
Study & Research object & Decision and conditions & Evaluation or claim boundary\\\midrule
Contract theory \cite{benveniste2018contracts} & Assumption--guarantee specifications and their composition & Component compatibility under stated assumptions & General design formalism; not a new algebra in this paper\\
Pant et al. \cite{pant2021} & Error/delay contracts for anytime estimation and control & Joint computation/control with robust MPC conditions & Physical aerial-platform and closed-loop control evaluation\\
Learn-then-test \cite{angelopoulos2025ltt} & Risk-controlled predictive configurations & Multiple tests admit configurations from calibration & Statistical risk control; used as an established selection principle\\
Runtime assurance \cite{mehmood2022simplex} & Safe autonomy with a fallback controller & Monitoring and switching under safety conditions & Control assurance; our evidence refusal supplies no physical safety theorem\\
This study & Correlated pose revisions and event evidence bound to their policy/source/installation identity & Pre-sensing commitment, compatible execution, source-scoped statement lifetime & Numerical, task and applicability endpoints tested separately against a complete procedural control\\\bottomrule
\end{tabularx}
\end{table}
